\subsection{欧氏球与球面}

对每个整数 \(p > 0\) ，令 \(\mathbf{U}_p\) 表示由一切点对 \((x, y)\) （\(\mathbf{R}^n\) 中的点对）组成的满足 \(d(x, y) < 1/p\) 的集；不等式 (3) 表明，诸集 \(\mathbf{U}_p\) 构成 \(\mathbf{R}^n\) 的一致结构的一个基本围域系（参见第 IX 章，§ 2）。

由这个事实以及不等式

\[
\left| d \left(\boldsymbol {x}, \boldsymbol {y}\right) - d \left(\boldsymbol {x} ^ {\prime}, \boldsymbol {y} ^ {\prime}\right) \right| \leqslant d \left(\boldsymbol {x}, \boldsymbol {x} ^ {\prime}\right) + d \left(\boldsymbol {y}, \boldsymbol {y} ^ {\prime}\right),
\]

——它是 (1) 的推论——我们推出 \(d(x, y)\) 在 \(\mathbf{R}^n \times \mathbf{R}^n\) 上一致连续；因而范数 \(||x|| = d(0, x)\) 在 \(\mathbf{R}^n\) 上一致连续。

定义 1. 给定一点 \(\mathbf{x}_0 \in \mathbb{R}^n\) 与一个实数 \(r > 0\) ，\(n\) 维的、以 \(\mathbf{x}_0\) 为中心、以 \(r\) 为半径的开（相应地，闭）欧氏球，是由所有点 \(\mathbf{x} \in \mathbb{R}^n\) ——满足 \(d(\mathbf{x}_0, \mathbf{x}) < r\) 的——组成的集{[}相应地，满足 \(d(\mathbf{x}_0, \mathbf{x}) \leqslant r\){]}；\(n - 1\) 维的、以 \(\mathbf{x}_0\) 为中心、以 \(r\) 为半径的欧氏球面，是由所有 \(\mathbf{x} \in \mathbb{R}^n\) ——满足 \(d(\mathbf{x}_0, \mathbf{x}) = r\) 的——组成的集。

当没有混淆的危险时，我们把“欧氏球”简称为“球”（相应地，把“欧氏球面”简称为“球面”）。当 n = 2 时，二维的球称为圆盘，一维的球面称为圆周。当 n = 1 时，以 \(x_{0}\) 为中心、以 r 为半径的开（相应地，闭）球是区间 \(]x_{0} - r, x_{0} + r[\) （相应地，\([x_{0} - r, x_{0} + r]\)）；以 \(x_{0}\) 为中心、以 r 为半径的球面是由这两个区间的两个端点 \(x_{0} - r, x_{0} + r\) 组成的集。

由上面所述，以 \(x_{0}\) 为中心的球（开的或闭的）（或者只是那些以 1/p 为半径的球，其中 p 取遍整数 \textgreater0 的集）构成点 \(x_{0}\) 的一个基本邻域系。

命题 1. \(\mathbf{R}^n\) 的每个开（相应地，闭）球是开集（相应地，紧集）。开球的闭包是同一中心、同一半径的闭球；闭球的内部是同一中心、同一半径的开球。

以 \(x_0\) 为中心、以 \(r\) 为半径的开（相应地，闭）球是区间 \(] - \infty, r[\) （相应地，\(] - \infty, r]\)）在连续函数 \(d(x_0, x)\) 下的逆象；因此它是开的（相应地，是闭的且有界的，从而是紧的）。若 \(d(x_0, x) = r\) ，且 \(y = x_0 + t(x - x_0)\) （ \(0 < t < 1\) ）是以 \(x_0\) 与 \(x\) 为端点的开线段上的一个点，则我们有 \(d(x_0, y) = tr < r\) ，而 \(d(x, y) = (1 - t)r\) 可以任意小；因此 \(x\) 属于以 \(x_0\) 为中心、以 \(r\) 为半径的开球的闭包。同样，若 \(z = x + t(x - x_0)\) （ \(t > 0\) ）是以 \(x\) 为端点、以 \(x - x_0\) 为方向向量的开射线上的一个点，则我们有

\[
d (x _ {0}, z) = (1 + t) r > r,
\]

而 \(d(\boldsymbol{x}, \boldsymbol{z}) = tr\) 可以任意小；因此 x 不是以 \(x_{0}\) 为中心、以 r 为半径的闭球的内点。

推论. 每个欧氏球面是紧集，并且是同一中心、同一半径的开球与闭球的边界。

映射 \(x \to \frac{1}{r} (x - x_0)\) 把以 \(x_0\) 为中心、以 \(r\) 为半径的球面（相应地，开球、闭球）变换成以 \(o\) 为中心、以 \(1\) 为半径的球面（相应地，开球、闭球）；这个球面记作 \(S_{n-1}\) ，称为 \(R^n\) 中的单位球面。同样，以 \(o\) 为中心、以 \(1\) 为半径的闭球记作 \(B_n\) ，称为 \(R^n\) 中的单位球。\((n - 1)\) 维球面（相应地，\(n\) 维闭球）的拓扑研究因此归结为 \(S_{n-1}\) （相应地，\(B_n\)）的拓扑研究。对开球，我们有下述命题：

命题 2. 每个 n 维开球同胚于 \(R^{n}\)。

因为映射 \(x \to \frac{x}{1 + ||x||}\) 在 \(\mathbf{R}^n\) 上连续，并且把 \(\mathbf{R}^n\) 映到以 \(0\) 为中心、以 \(1\) 为半径的开球上；此外，由 \(y = \frac{x}{1 + ||x||}\) 我们推出 \(x = \frac{y}{1 - ||y||}\) ，所以这个映射是双射的且双连续的。

令 \(\mathbf{R}_n^*\) 表示 o 在 \(\mathbf{R}^n\) 中的补集。

命题 3. 空间 \(\mathbf{R}_n^*\) 同胚于 \(\mathbf{S}_{n-1}\) 与空间 \(\mathbf{R}_+^*\) （由实数 \(>0\) 组成的）的乘积。

因为每一点 \(x \neq 0\) 都可唯一地写成形式 \(tz\) ，其中 \(t > 0\) 且 \(||z|| = 1\) ，因为由 \(x = tz\) 可推出 \(t = ||x||\) 与 \(z = x / ||x||\) 。由于 \(tz\) 在乘积空间 \(\mathbf{R} \times \mathbf{R}^n\) 上连续，从而更在 \(\mathbf{R}_+^* \times \mathbf{S}_{n-1}\) 上连续，并且 \(||x||\) 与 \(\frac{1}{||x||}\) 在 \(\mathbf{R}_n^*\) 上连续，命题得证。

映射 \(x \to x / ||x||\) 称为 \(\mathbf{R}_n^*\) 到 \(\mathbf{S}_{n-1}\) 上的中心投影。同样地定义一点 \(a\) 的补集到以 \(a\) 为中心的球面上的中心投影。

推论 1. 球面 \(S_{n-1}\) 同胚于 \(R_{n}^{*}\) 关于下述等价关系的商空间：该等价关系的类是以 o 为端点的开射线。

这些类也可以定义为 \(\{0\}\) 以外的、比值 \textgreater0 的位似群的非传递类。

推论 2. 空间 \(\mathbf{R}_n^*\) 同胚于 \(\mathbf{R} \times \mathbf{S}_{n-1}\)。

因为 \(\mathbf{R}_{+}^{*} = ]0, +\infty[\) 同胚于 \(\mathbf{R}\)（第 IV 章，§ 4，第 1 小节，命题 1）。

注. 这些命题并非欧氏球所特有，而可以推广到 \(R^{n}\) 中 o 的整整一类紧邻域（见习题 12）。

集 \(S_{n-1}\) 与 \(B_{n}\) 显然在一切正交变换下不变。若 V 是 \(R^{n}\) 中一个 p 维向量子空间，则存在把 V 变换成 p 维坐标簇的正交变换；因此 \(V \cap S_{n-1}\) （相应地，\(V \cap B_{n}\)）同胚于 \(S_{p-1}\) （相应地，\(B_{p}\)）。

